\section{Results}
\label{sec:results}

\subsection{Cross-phase Summary}
% Pull from experiments/PHASE_REPORT_FULL.md cross-phase table.
\begin{table}[t]
\centering
\caption{Cross-phase Mean F1 macro and Mean Accuracy. Best per phase in bold.}
\label{tab:cross_phase}
\small
\begin{tabular}{@{}llll@{}}
\toprule
Phase & Method                & Mean F1 macro       & Mean Acc \\
\midrule
1     & Baseline (no CBAM)    & 0.420               & 81.4\,\% \\
2     & CBAM + sqrt CW        & 0.509               & 68.6\,\% \\
2     & \textbf{Hybrid (E7)}  & \textbf{0.516}      & 70.0\,\% \\
\midrule
3     & Hybrid 0-shot         & 0.362               & --       \\
3     & \textbf{Naked BMC 0-shot} & \textbf{0.394}  & --       \\
\midrule
4     & Vanilla SPIDER        & 0.610               & 77.1\,\% \\
4     & CBAM SPIDER           & 0.597               & 75.2\,\% \\
4     & \textbf{Hybrid frozen} & \textbf{0.623}     & 79.7\,\% \\
4     & Hybrid unfreeze       & 0.622               & 80.4\,\% \\
\bottomrule
\end{tabular}
\end{table}

\subsection{Phase 1+2 — RSNA Severe-class Recovery}
\begin{table}[t]
\centering
\caption{Severe-class F1 on RSNA validation (per condition).}
\label{tab:severe_f1}
\small
\begin{tabular}{@{}lccc@{}}
\toprule
Method            & Spinal Severe & L. Foram. Severe & R. Foram. Severe \\
\midrule
Baseline          & 0.457         & \textbf{0.000}   & \textbf{0.000} \\
CBAM + sqrt CW    & \textbf{0.623}& 0.180            & 0.197 \\
Hybrid (E7)       & 0.500         & \textbf{0.276}   & \textbf{0.283} \\
\bottomrule
\end{tabular}
\end{table}

The CBAM spatial gate recovers Foraminal Severe F1 from $0.000$
(baseline collapse) to $0.180/0.197$. Adding the frozen BiomedCLIP
branch (Hybrid) further improves to $0.276/0.283$, showing the
language-aligned semantic features supply complementary signal beyond
3D attention.

\subsection{Phase 3 — Zero-shot SPIDER}
\begin{figure}[t]
  \centering
  \includegraphics[width=\columnwidth]{cross_phase_comparison.png}
  \caption{Mean F1 macro across all phases; red border indicates best
           in phase. Selective transfer is evident in Phase 3, where
           the Hybrid wins on disc-related labels but a 2D-only Naked
           BiomedCLIP wins on vertebra/endplate labels.}
  \label{fig:cross_phase}
\end{figure}

\begin{table}[t]
\centering
\caption{Per-disease F1 macro on SPIDER zero-shot evaluation
         (no SPIDER training used). $\Delta$ = Hybrid $-$ Naked.}
\label{tab:spider_zeroshot}
\small
\begin{tabular}{@{}llccr@{}}
\toprule
Disease            & Tier   & Hybrid          & Naked BMC         & $\Delta$ \\
\midrule
Disc\_narrowing    & easy   & \textbf{0.586}  & 0.403             & $+0.183$ \\
Disc\_bulging      & medium & \textbf{0.613}  & 0.363             & $+0.250$ \\
Disc\_herniation   & medium & \textbf{0.563}  & 0.532             & $+0.031$ \\
Pfirrman           & hard   & \textbf{0.155}  & 0.152             & $+0.003$ \\
\midrule
UP\_endplate       & hard   & 0.467           & \textbf{0.569}    & $-0.102$ \\
LOW\_endplate      & hard   & 0.468           & \textbf{0.558}    & $-0.090$ \\
Spondylolisthesis  & medium & 0.028           & \textbf{0.500}    & $-0.472$ \\
Modic              & hard   & 0.018           & \textbf{0.071}    & $-0.053$ \\
\midrule
\textbf{Mean (8)}  &        & 0.362           & \textbf{0.394}    & $-0.032$ \\
\bottomrule
\end{tabular}
\end{table}

\begin{figure}[t]
  \centering
  \includegraphics[width=\columnwidth]{ablation_hybrid_vs_naked.png}
  \caption{Phase 3 ablation: per-disease zero-shot F1 macro.
           Hybrid wins on disc-related labels (top), Naked BiomedCLIP
           wins on vertebra/endplate labels (bottom) — supporting the
           complementarity argument.}
  \label{fig:ablation}
\end{figure}

The split in Tab.~\ref{tab:spider_zeroshot} and Fig.~\ref{fig:ablation}
is the central evidence for selective transfer: gains on disc-related
labels (where 3D volumetric context dominates) trade against losses on
vertebra/endplate labels (where 2D appearance is sufficient and the
RSNA-trained projection MLP over-specialises). We interpret this as a
characterisation of \emph{when} the hybrid integration is preferred to
a frozen 2D-only baseline, rather than a failure mode.

\subsection{Phase 4 — SPIDER Transfer Learning}
\begin{table}[t]
\centering
\caption{Per-condition F1 macro on SPIDER supervised transfer.
         Best per row in bold. Hybrid frozen achieves the best mean
         with ${\sim}5$M trainable parameters vs $22$M for Vanilla.}
\label{tab:spider_transfer}
\small
\setlength{\tabcolsep}{4pt}
\begin{tabular}{@{}lcccc@{}}
\toprule
Condition           & Vanilla & CBAM   & Hybrid fz       & Hybrid unfz     \\
\midrule
Pfirrmann (5-cls)   & 0.594   & 0.535  & \textbf{0.626}  & 0.580           \\
Modic (4-cls)       & 0.363   & 0.358  & 0.373           & \textbf{0.387}  \\
Disc Narrow.\ (2)   & 0.851   & 0.862  & 0.865           & \textbf{0.878}  \\
Spondylolist.\ (2)  & 0.633   & 0.634  & 0.627           & \textbf{0.643}  \\
\midrule
Mean F1 macro       & 0.610   & 0.597  & \textbf{0.623}  & 0.622           \\
Mean Acc            & 77.1\,\% & 75.2\,\% & 79.7\,\%     & \textbf{80.4\,\%} \\
Trainable params    & 22M     & 22M    & \textbf{5M}     & 95M             \\
\bottomrule
\end{tabular}
\end{table}

With $\sim 5$M trainable parameters, Hybrid frozen reaches Mean F1
macro $0.623$, exceeding both fully fine-tuned Vanilla ($0.610$) and
CBAM ($0.597$) baselines while training in $\sim 5$ minutes per epoch
on a single GPU. Full backbone fine-tuning (Hybrid unfreeze) does not
improve mean F1 — gains on three conditions are offset by a Pfirrmann
regression — supporting frozen as the better default
(Tab.~\ref{tab:spider_transfer}).
